\section{Weighted tests and the universal AME bound}
\label{app:marginal-verification}

This appendix refines Proposition~\ref{prop:complementary-marginal-tests}.
The first result quantifies the effect of selecting or reweighting tests.
The second removes the stabilizer assumption from the uniform gap and
weighted lower bound.  Neither refinement is needed for the rigidity
proofs or the observable criterion in
Corollary~\ref{cor:marginal-certified-rounding}.

For a perfectly complete verification operator \(\Omega\), write
\[
 \nu(\Omega)=\min_{v\perp\psi,\ \|v\|=1}
                  \langle v|(I-\Omega)|v\rangle.
\]
A permitted binary effect is \(E_S\otimes I_{S^c}\), with
\(0\leq E_S\leq I\), \(|S|\leq m+1\), and
\((E_S\otimes I)\psi=\psi\).  Randomization chooses among these fixed
supports before testing the copy.  Positivity implies that \(I-E_S\)
annihilates the marginal support.  Consequently \(E_S\otimes I\geq\Pi_S\)
when \(|S|=m+1\); for \(|S|\leq m\), full marginal rank forces
\(E_S=I\).  This is why half-plus-one is the first informative support size.

\begin{proposition}[Weights and number of tests]
\label{prop:weighted-marginal-tests}
For a stabilizer \(\AME(2m,q)\), \(m\geq2\), assign probabilities
\(w_S\) to the complementary library \(\mathcal S\), including zero
weights, and order them as \(w_{(1)}\leq\cdots\leq w_{(2m)}\).  Then
\begin{equation}
 \nu\!\left(\sum_{S\in\mathcal S}w_S\Pi_S\right)
       =\sum_{a=1}^{m+1}w_{(a)}.
 \label{eq:weighted-marginal-gap}
\end{equation}
Among all permitted protocols with at most \(s\) binary tests,
\begin{equation}
 \nu_{\max}(s)=
 \begin{cases}
 0,&1\leq s<m,\\
 1-(m-1)/s,&m\leq s\leq2m,\\
 (m+1)/(2m),&s\geq2m.
 \end{cases}
 \label{eq:marginal-test-tradeoff}
\end{equation}
Any \(s\) distinct library tests sampled uniformly attain the middle
branch; the full library attains the last.
\end{proposition}

\begin{proof}
The character argument in Proposition~\ref{prop:complementary-marginal-tests}
shows that every nontrivial character fails at least \(m+1\) tests, and
that any prescribed \(m-1\) tests are exactly the pass set of some
nontrivial character.  Let it pass the \(m-1\) heaviest weights to obtain
\eqref{eq:weighted-marginal-gap}.

For an arbitrary collection of permitted supports, each contributes at
most \(2e\) independent character constraints.  Fewer than \(m\) tests
therefore have a common nontrivial passing character.  If there are
\(t\geq m\) positive weights, a nontrivial character passes the support
projectors of the \(m-1\) heaviest tests, and hence their effects.  Its
rejection probability is at most \(1-(m-1)/t\).  Combine this with the
one-site-error cap \((m+1)/(2m)\) proved in the body.  Uniform sampling
from the library attains both bounds in their stated ranges.
\end{proof}

In particular, uniform weights uniquely maximize the gap within the full
library.  If \(f\leq m\) tests become unavailable, the remaining tests
still give gap \((m+1-f)/(2m-f)\) when sampled uniformly.  Here the
parties remain available; only the list of tests is reduced.
The general minimum-test gap bound and common-eigenbasis optimization
precede this calculation \cite[Proposition~1 and Section~V.B]{DangniamHanZhu2020}.

\begin{proposition}[Arbitrary AME targets]
\label{prop:universal-marginal-verification}
Let \(\psi\) be any \(\AME(2m,q)\), with integers \(q\geq2,m\geq2\).
For the same complementary library and arbitrary probability weights,
\begin{equation}
 I-\sum_{S\in\mathcal S}w_S\Pi_S
 \geq\left(\sum_{a=1}^{m+1}w_{(a)}\right)
       (I-|\psi\rangle\langle\psi|).
 \label{eq:universal-weighted-gap}
\end{equation}
Any \(m\) library marginals determine \(\psi\) among all density
operators.  The uniform full library has gap exactly \((m+1)/(2m)\),
optimal among permitted tests.  Exact weighted attainment and the
finite-test upper bound above retain their stabilizer hypothesis.
\end{proposition}

\begin{proof}
We replace stabilizer characters by errors on a half-set.  Choose on each
site a unitary operator basis orthonormal for normalized Hilbert--Schmidt
inner product, with identity distinguished.  Maximal mixing on \(B\)
makes the \(q^{2m}\) vectors \(W_B\psi\) an orthonormal basis.  The
projector \(\Pi_{C\cup\{i\}}\) has rank \(q^{2m-2}\); it fixes the
\(q^{2m-2}\) basis vectors whose error at \(i\) is identity, since those
errors act on its complement.  These vectors therefore span its range.
Each star is consequently a commuting family, testing one error coordinate
at a time, with product \(|\psi\rangle\langle\psi|\).
Equivalently, a half-unitary taking \(\psi\) to \(m\) Bell pairs
conjugates these projectors to the pair projectors.  This is the
virtual-factor picture of parent constraints
\cite[Sections~IV--V]{JohnsonTicozziViola2017}
\cite[Sections~2 and~5]{KaruvadeJohnsonTicozziViola2018}.

The two stars need not commute with each other.  Orthogonality of their
low-rejection spaces replaces a common eigenbasis in the following argument.

Select \(r\) tests from this star and \(t\) from the other, where
\(s=r+t\geq m\), and put \(d=s-m\).  On \(\psi^\perp\), let \(A,D\)
be the respective sums of rejection projectors and let
\(E_k=\mathbf1_{A\leq k}\), \(F_l=\mathbf1_{D\leq l}\).
The range of \(E_k\) is spanned by nonidentity errors on \(B\) with at
most \(k\) errors on the \(r\) tested coordinates, hence total weight at
most \(m-r+k\).  The analogous bound for \(F_l\) is \(m-t+l\).
If \(k+l\leq d\), the combined support of two such opposite-half errors
has size at most \(m\).  Maximal mixing makes their inner product zero:
the errors are on disjoint halves and each vector has a nonidentity factor
of trace zero.  Thus \(E_kF_l=0\).

In particular \(E_k+F_{d-k}\leq I\) for \(0\leq k\leq d\).
The spectra of \(A,D\) are nonnegative integers, so spectral calculus gives
\[
 \sum_{k=0}^d E_k=(d+1)I-\min\{A,(d+1)I\}.
\]
Summing the projection inequalities gives
\[
 \min\{A,(d+1)I\}+\min\{D,(d+1)I\}\geq(d+1)I.
\]
Each truncated operator is bounded above by its untruncated counterpart,
so \(A+D\geq(d+1)I\); no joint functional calculus is used.
Restoring the target ray, every selected family of \(s\geq m\) tests
therefore satisfies the \emph{selected-family bound}
\begin{equation}
 \sum_{S\text{ selected}}(I-\Pi_S)
       \geq(s-m+1)(I-|\psi\rangle\langle\psi|).
 \label{eq:selected-family-verification}
\end{equation}
Order the weights, write \(w_0=0\), and expand the weighted sum as
\[
 \sum_{i=1}^{2m}w_i(I-\Pi_i)
  =\sum_{k=1}^{2m}(w_k-w_{k-1})
                    \sum_{i=k}^{2m}(I-\Pi_i).
\]
Apply the selected-family bound to tails of size at least \(m\), and
positivity to the rest.  The coefficient telescopes to
\(\sum_{i=1}^{m+1}w_i\), proving \eqref{eq:universal-weighted-gap}.
A positive gap for any \(m\) selected tests gives marginal determination.
Uniform full sampling attains the one-site-error upper bound from the body,
whose proof uses only one-party maximal mixing.
\end{proof}
